\section{Limits and threats to validity}
\label{sec:limits}

\paragraph{Physical model.} The case-defined reference house assumes a lossless tank, 80\% heat
recovery, constant daily hot-water demand, flat electricity price and abstract \qty{1}{kW}
servers. The smaller four-L4 variant uses host-power estimates and GPU TDP rather than measured
wall power or a quoted liquid-cooled build. Its calculated heat fits the constant hot-water
base load; actual draw timing, delivery temperature, thermal losses and water hygiene require
engineering validation. Neither normalised 1~kW L4 profile nor the candidate four-card host is
an installed bill of materials.

\paragraph{Work and markets.} The reference-house maximum assumes unlimited paid work. The
mixed-queue run uses a medical-dictation arrival fixture for technical validation, not the
proposed public-document LLM market; sensitive health records are not placed in homes. Its
0.35 filler capacity cap is an analogy, not observed sales. The laptop Apertus measurement is
CPU-package energy, not L4 throughput or whole-node power. The research simulation was measured
on a CPU laptop; a billable GPU-compatible filler has not been demonstrated. Competitor list
prices and retail GPU rates are buyer prices, not DigitalHeat contracts or net receipts. GPU wear
and purchase prices are sensitivity inputs. Cooling integration, idle energy, connectivity,
maintenance, insurance, financing, tax bases and sales costs remain incompletely priced. The
consumer-GPU cost case is excluded from the proposed deployment pending licence clarification.

\paragraph{Network inference.} The 33-site locations and sizes and Type~1 arrival shares are
synthetic. Unless a forecast-error sweep is explicitly stated, planning uses perfect heat
foresight while knowing only arrived jobs. Uniform transfer delay assumes data staged at arrival;
transfer cost and distance-dependent delay are absent. The experiments therefore compare
algorithms under controlled inputs rather than predict an operating Swiss network.
